\documentclass[12pt]{article}
\usepackage{amsmath, amsthm, amssymb}
\usepackage[margin=1.25in]{geometry}
\usepackage{microtype}
\usepackage{booktabs}
\usepackage{hyperref}

\newtheorem{recommendation}{Recommendation}

\title{Review of \texttt{APS\_repair.pdf}:\\
Remaining Errors and Recommendations}
\author{Review against \texttt{ch24\_revisions\_claude.pdf}}
\date{\today}

\begin{document}
\maketitle

\begin{abstract}
The document \texttt{APS\_repair.pdf} successfully incorporates the principal
logical corrections recommended in \texttt{ch24\_revisions\_claude.pdf}:
Lemma~2, Definitions~7 and~8, the monotonicity proof, Theorem~1, the
compactness proposition, and the two-value corollary are all correctly
stated and correctly ordered.  Three expository residues of the old
four-tuple formulation remain, however, and one algorithmic passage
should either be rewritten or given an explicit citation to the corollary
that licenses it.  This note identifies each issue and states the required
correction.
\end{abstract}

\section{Background}

The core repair identified in \texttt{ch24\_revisions\_claude.pdf} is the
replacement of every admissible \emph{four-tuple} $(x,y,w_1,w_2)$ by an
admissible \emph{triple} $(x,y,w(\cdot))$, where $w(\cdot):Y\to W$ is a
function, throughout the argument that establishes compactness of $V$.
The two-value characterization---using scalar continuation values
$(v_1,v_2)\in V\times V$---is then recovered as a \emph{corollary} once
compactness is in hand.  The logical structure of the revised chapter is:

\begin{center}
Lemma~2 (triples) $\;\Rightarrow\;$
Definitions~7--8 (triples) $\;\Rightarrow\;$
Theorem~1 $\;\Rightarrow\;$
Proposition~2 (compactness) $\;\Rightarrow\;$
Corollary (two-value form).
\end{center}

\section{Errors That Remain in \texttt{APS\_repair.pdf}}

\subsection{Error~1: Four-tuple language in the paragraph following
Definition~7 (p.\,1141)}

After stating Definition~7 correctly in triple form, the immediately
following paragraph reverts to the old language.  The text reads:

\begin{quote}
``Notice that when $W\subset V$, the admissible 4-tuple
$(x,y,w_1,w_2)$ determines an SPE with strategy profile
$\sigma_1=(x,y)$, $\sigma|_{(x,y)}=\sigma^1$,
$\sigma|_{(\nu,\eta)}=\sigma^2$ for $(\nu,\eta)\neq(x,y)$
where $\sigma^1$ is a continuation strategy that yields value
$w_1=V_g(\sigma^1)$ and $\sigma^2$ is a strategy that yields
continuation value $w_2=V_g(\sigma^2)$.  The value of the SPE is
$V_g(\sigma)=w=(1-\delta)r(x,y)+\delta w_1$.''
\end{quote}

This paragraph introduces scalar continuation values $w_1$ and $w_2$
under the label ``admissible 4-tuple,'' contradicting the corrected
Definition~7 that precedes it.

\begin{recommendation}
Replace the paragraph with the following:

\smallskip
\noindent
Notice that when $W\subset V$, an admissible triple $(x,y,w(\cdot))$
determines an SPE with strategy profile
\[
  \sigma_1 = (x,y), \qquad
  \sigma|_{(x,\eta)} = \tilde\sigma(\eta) \quad \forall\,\eta\in Y,
\]
where $\tilde\sigma(\eta)$ is any SPE with value $w(\eta)\in V$.  The
value of the SPE is
\[
  V_g(\sigma) = (1-\delta)\,r(x,y)+\delta\,w(y).
\]
\end{recommendation}

\subsection{Error~2: Two occurrences of ``4-tuples'' in Section~25.8.1
(p.\,1141--1142)}

Section~25.8.1 (``Basic idea of dynamic programming squared'') contains
two sentences that still use four-tuple language:

\begin{enumerate}
\item ``APS define a mapping $B(W)$ that, by considering only admissible
      \textbf{4-tuples}, maps a \emph{set} of values $W$ tomorrow into a
      new \emph{set} $B(W)$ of values today.''
\item ``Thus, APS use admissible \textbf{4-tuples} to map candidate
      continuation values tomorrow into new candidate values today.''
\end{enumerate}

\begin{recommendation}
Replace both occurrences of ``4-tuples'' with ``triples,'' so that the
two sentences read:

\begin{enumerate}
\item ``APS define a mapping $B(W)$ that, by considering only admissible
      \textbf{triples}, maps a set of values $W$ tomorrow into a new set
      $B(W)$ of values today.''
\item ``Thus, APS use admissible \textbf{triples} to map candidate
      continuation values tomorrow into new candidate values today.''
\end{enumerate}
\end{recommendation}

\subsection{Error~3: The iteration algorithm in Section~25.9 uses the
two-value form without citing the corollary (p.\,1145--1146)}

The algorithm for computing the boundaries of $B(W_0)$ states the
problem for the worst value $\underline{w}_1$ as a minimization over
$(x,y)\in C$ and $(w_1,w_2)\in[\underline{w}_0,\overline{w}_0]^2$:
\[
\underline{w}_1 =
\min_{\substack{(x,y)\in C\\(w_1,w_2)\in[\underline{w}_0,\overline{w}_0]^2}}
(1-\delta)\,r(x,y)+\delta w_1
\]
subject to
\[
(1-\delta)\,r(x,y)+\delta w_1 \;\ge\;
(1-\delta)\,r(x,\eta)+\delta w_2
\quad\forall\,\eta\in Y.
\]
This is the four-tuple formulation.  It is mathematically correct
\emph{given} the two-value Corollary, but in the revised chapter the
Corollary appears \emph{before} this algorithm on p.\,1144--1145.
The algorithm therefore either (a)~should be rewritten in triple form
for full consistency, or (b)~must include an explicit forward reference to
the Corollary that licenses the reduction to two scalar values.  As
currently written it does neither, leaving the reader without justification
for the move from $w(\cdot)$ back to $(w_1,w_2)$.

\begin{recommendation}
Add a sentence immediately before the two optimization problems:

\smallskip
\noindent
``By the two-value Corollary above, it suffices to search over pairs
$(w_1,w_2)\in W\times W$ rather than full functions $w(\cdot):Y\to W$;
the boundary problems therefore reduce to:''

\smallskip
\noindent
Alternatively, rewrite the optimization entirely in triple form:
\[
\underline{w}_1 =
\min_{(x,y,w(\cdot))\text{ admissible w.r.t.\,}[\underline{w}_0,\overline{w}_0]}
(1-\delta)\,r(x,y)+\delta\,w(y),
\]
and note that, by the Corollary, the minimum is attained at a function
$w(\cdot)$ that takes at most two values.
\end{recommendation}

\section{Summary Table}

\begin{center}
\begin{tabular}{llp{6cm}}
\toprule
\textbf{Location} & \textbf{Issue} & \textbf{Required action} \\
\midrule
p.\,1141, paragraph after Def.\,7
  & ``admissible 4-tuple $(x,y,w_1,w_2)$''
  & Replace with triple $(x,y,w(\cdot))$ description \\[4pt]
p.\,1141--42, \S25.8.1 (twice)
  & ``admissible 4-tuples'' (two occurrences)
  & Replace both with ``admissible triples'' \\[4pt]
p.\,1145--46, algorithm in \S25.9
  & Minimization over $(w_1,w_2)\in W^2$ without justification
  & Add citation to two-value Corollary, or rewrite in triple form \\
\bottomrule
\end{tabular}
\end{center}

\section{Items Correctly Implemented}

For completeness, the following elements of
\texttt{ch24\_revisions\_claude.pdf} are correctly incorporated:

\begin{itemize}
\item \textbf{Lemma~2} (p.\,1140): states and proves the necessary and
  sufficient condition using a function $v(\cdot):Y\to V$.
\item \textbf{Definition~7} (p.\,1141): correctly defines admissible
  \emph{triples} $(x,y,w(\cdot))$.
\item \textbf{Definition~8} (p.\,1142): $B(W)$ defined via admissible
  triples.
\item \textbf{Monotonicity proof} (p.\,1142--43): uses the
  triple/function formulation throughout.
\item \textbf{Theorem~1 proof} (p.\,1143--44): constructs sequences of
  functions $w_1(\cdot),w_2(\cdot),\ldots$ correctly.
\item \textbf{Proposition~2} (compactness, p.\,1144): present, correctly
  placed, and correctly proved.
\item \textbf{Corollary} (two-value characterization, p.\,1144--45):
  correctly positioned \emph{after} compactness, with a correct
  necessity proof using $v_2=\inf_{\eta\in Y}v(\eta)\in V$.
\end{itemize}

\end{document}
